\subsection{Step: amplify an interaction into learned routing}
A running count needs one routing decision: read the displayed count and the current letter.
We give one block-$2$ head, the gate head, a segment mask under which each count position attends to itself and all $n$ letters.
A relative bias $\beta$ determines which letter is selected; the current letter has offset $n$.
This head supports a second exact step solution besides the lookup executor of \cref{prop:represent}.

\begin{proposition}[Gate executor]
\label{prop:gate}
For counting modulo~$K$ with the gate head, there are finite relative biases and a readout, giving $\theta^\star_{\mathrm{gate}}$, at
which the attention outside the count's own position and the current letter changes every margin by less than $\tfrac12$, the
two are balanced, and every prediction $s_t$, $t\ge2$, is correct with margin at least $\tfrac12$.
\end{proposition}

\paragraph{A weak pathway at initialization.}
Write $\bm x_t=\bm e_{s_{t-1}}$ for the displayed count and $m_t=(-1)^{b_t}$ for the letter sign.
The structured seed in \cref{app:routing} reduces the state logits to
\[
 \bm z_t=\initscale U\bm x_t+g\,(1-M_t)\,W\bm x_t,
 \qquad |a|=\frac{\initscale^4}{\sqrt{d_{\mathrm h}}},\quad |g|=\initscale^2.
\]
Here $U$ and $W$ are $K\times K$ readout blocks, $M_t$ is the gate head's total attention mass on letters, so $1-M_t$ is the weight
the count keeps on its own position, and $a$ is the signed interaction strength in each letter score $am_j+\beta_{i_t-\ell_j}$; the
own-count score is $\beta_0$. The key reads the letter and the value of the count position carries the displayed count, so the gate
output $g(1-M_t)\bm x_t$ is a count$\times$letter product. Both readout blocks and all relative biases start at zero:
\[
 U(0)=W(0)=0,\qquad \beta(0)=0.
\]
The initialization supplies a weak count--letter pathway. At zero bias every letter is treated symmetrically;
the seed supplies neither a selected current letter nor a transition classifier.
The shared-presence control confirms why the pathway matters: if the key reads a direction both letters share,
$M_t$ does not depend on the letters and no bias or readout can outperform the count-only prior (\cref{lem:wrong-seed}).

\paragraph{Four stages.} Each stage is gradient flow on one parameter block, with all other parameters frozen:
\[
 \boxed{\begin{array}{cll}
 \mathrm A &U&\text{fit the count-only prior}\\[1mm]
 \downarrow&&\\[1mm]
 \mathrm B &W&\text{give the local interaction nonzero amplitude}\\[1mm]
 \downarrow&&\\[1mm]
 \mathrm C &\beta&\text{route to the own count and current letter}\\[1mm]
 \downarrow&&\\[1mm]
 \mathrm D &(U,W)&\text{fit the transition and drive loss to zero}
 \end{array}}
\]
Stage A fits the state-only predictor $\bm q_s$. Stage B stops at a small gate-readout level;
Stage C stops when attention leakage outside the two useful sites is at most $\zeta$; Stage D keeps $\beta$ frozen.
\Cref{app:routing} specifies the stopping levels and proves the stage bounds.

\paragraph{Why the stages work.} Write $s=s_{t-1}$ and $\partial_b\Psi=\Psi|_{b_t=1}-\Psi|_{b_t=0}$ for the response of a feature to the
current letter. At the count-only prior, every bounded feature $\Psi$ of the displayed count and letters obeys
\begin{equation}
 \E\bigl[(\bm e_{s_t}-\bm q_{s})\Psi\bigr]=p(1-p)\,\E\bigl[(\bm e_{s+1}-\bm e_{s})\,\partial_b\Psi\bigr],
 \qquad p(1-p)=\frac{1-\mixrate^2}{2-2\cos(2\pi/K)}
 \label{eq:identity}
\end{equation}
(\cref{lem:R-identity}). Everything already explained by the displayed state cancels. The gradient extracts only how a feature
responds to the current letter, along the move $\bm e_{s+1}-\bm e_s$ of the rule signal, with the local innovation paid once.
For the gate output, \cref{eq:identity} drives $W$ along the cyclic shift minus the identity, the readout that maps each displayed count to its move.
Stage B gives this interaction a nonzero readout coefficient so that it can drive the bias gradient.

For the common useful-site bias $\bar\beta=(\beta_n+\beta_0)/2$ and imbalance $\Delta=\beta_n-\beta_0$, Stage C establishes
\[
 \dot{\bar\beta}>0,\qquad
 \dot\beta_\delta\le0\quad(\delta\notin\{0,n\}),\qquad
 |\Delta|\le1.
\]
The useful sites gain mass, irrelevant offsets are suppressed, and the two useful sites remain balanced.
An invariant-region argument bounds the leakage hitting time by a polynomial (\cref{lem:R-field,lem:R-C}).
After routing, $1-M_t$ is a small perturbation of a feature that takes one value per letter, so a linear readout of the gate
output gives $\bm e_{s_{t-1}}+b_t(\bm e_{s_{t-1}+1}-\bm e_{s_{t-1}})=\bm e_{s_t}$; the final readout loss is convex and has a uniform
positive-margin separator. This gives the positive optimization result.

For $\zeta>0$, define the \emph{gate region} $\mathcal G_\zeta$ by leakage at most $\zeta$ outside
$\{$own position, current letter$\}$, $|\beta_n-\beta_0|\le1$, and positive margin on every prediction of $\Lstep$.

\begin{theorem}[Step training learns its routing]
\label{thm:routing}
Consider counting modulo~$K$ with a fixed query of probability $p\in(0,1)$ and the two-block Transformer with the gate
head, started from the small structured initialization of \cref{app:routing}, which supplies a weak count$\times$letter
pathway but zero readout and relative bias. There is $\initscale_\star$ with $1/\initscale_\star\le C\sqrt w$ such that, for
$0<\initscale\le\initscale_\star$, staged block-coordinate gradient flow on $\Lstep$ enters $\mathcal G_\zeta$ by a time
$T\le\mathrm{poly}(n,K,d_{\mathrm h},\initscale^{-1},(p(1-p))^{-1})$, stays in it at every later time, and drives $\Lstep$ to
zero.
\end{theorem}
